\documentclass[11pt,a4paper]{article}
\usepackage[margin=25mm]{geometry}
\usepackage{amsmath,amssymb,amsthm}
\setlength{\emergencystretch}{2em}
\newtheorem{lemma}{Lemma}
\newcommand{\gen}[1]{\langle#1\rangle}
\newcommand{\Alt}{\operatorname{Alt}}
\title{The first five omission colours at a critical six-point head}
\author{Human Compression project: bounded L044 continuation}
\date{}
\begin{document}\maketitle
We retain the full standing hypotheses and explicit smaller-degree
alternating induction of the accepted critical-head theorem.
Thus $G\le A_n$ is transitive with connected string C-marking and
all maximal omissions intransitive, $n\ge12$, and the actual
first three generators form an extendable six-point Table~2 head.
Assume rank failure. That theorem gives $n=2r\ge18$ and, in the
literal accepted orientation,
\[
\begin{split}
 D&=\{1,\ldots,6\},\quad E=\Omega\setminus D,\quad Z=E\cup\{6\},\qquad
 G_{\ge3}=\Alt(Z),\\
 \rho_0&=(1\ 2)(3\ 4),\quad \rho_1=(2\ 3)(4\ 5),\\
 \rho_2&=(1\ 2)(5\ 6)\text{ or }(3\ 4)(5\ 6),\qquad
 \tau=\rho_3=(6\ b)\delta,\quad b\in E.
\end{split}
\]
Here $\delta$ is odd and supported on $E\setminus\{b\}$.
Put $\sigma=\rho_4$, $K=G_{\ge5}$ and $B=G_{\ge4}$.
The group $K$ fixes $D\cup\{b\}$ and centralizes $\delta$;
$\sigma$ fixes $D$ pointwise. No transitivity or natural alternating
identification is assumed for $B$.

\begin{lemma}[Colour four is never a unique crossing]
The actual omission $G_4$ has at least three orbits. Consequently
colour $4$ has at least two ambient-eligible edges.
\end{lemma}
\begin{proof}
Transitivity of $G_{\ge3}$ gives $x=b^\sigma\ne b$;
otherwise $\gen{\tau,\sigma,K}$ would preserve $\{6,b\}$.
The complete actual pure-head overlap argument gives $\tau x=x$:
conjugating the head first by $\tau$ and then by $\sigma$ supplies
pure $\Alt(D\cup\{b,x\})$ in $G_{\le4}$. If $\tau x\ne x$,
another overlap supplies pure $A_4$ in
$G_{\le4}\cap G_{\ge3}=\gen{\tau,\sigma}$, contrary to its
dihedral structure. In particular $\delta x=x$.

Let $P=G_{\le3}$. It preserves $W=D\cup\{b\}$ and is transitive
there; its exterior action is exactly $\gen\delta$.
The group $K$ fixes $W$ pointwise. Thus $W$ is one actual
$G_4=\gen{P,K}$ orbit, and the exterior orbits are those of
$\gen{\delta,K}$. Since $\delta$ commutes with $K$ and fixes $x$,
the exterior orbit of $x$ is exactly $x^K$ and is pointwise fixed
by $\delta$. The odd involution $\delta$ is nonidentity and fixes
all of $W$ as well, so some moved point belongs to a third orbit.

Adding the edges of $\rho_4$ to the omission graph must connect
all its orbits, by transitivity of $G$. A single eligible edge
can connect only two of them. Hence there are at least two.
This argument concerns full ambient omission orbits, not a local
omission in a projected parabolic.
\end{proof}

\begin{lemma}[The literal lower-colour routes]
Colours $0$ and $2$ have unique eligible edges; colour $1$ has
two. Their exact omission orbits are
\[
\begin{array}{c|c|c}
 &\text{row 2 head}&\text{row 3 head}\\ \hline
G_0&\{1,2,3\},\ \Omega\setminus\{1,2,3\}
    &\{1\},\ \Omega\setminus\{1\}\\
G_1&\multicolumn{2}{c}{\{1,2\},\ \{3,4\},\ E\cup\{5,6\}}\\
G_2&\multicolumn{2}{c}{\{1,2,3,4,5\},\ Z}.
\end{array}
\]
At colour $0$, the row~2 head has right-shared labels exactly
$1,2$, while the row~3 head has none. The colour-$2$ cut is
unshared. Colour $3$ has a unique eligible edge if and only if
$B$ is transitive on $E$; if so, its two sides are the actual
supported groups $A_5$ on $D$ and $B$ on $E$, and it is unshared.
\end{lemma}
\begin{proof}
The full tail $G_{\ge3}=\Alt(Z)$ is transitive on $Z$ and fixes
the first five points. Add the displayed remaining literal head
edges after each omission; they give exactly the listed components.
For $G_0$, the eligible edge is $\{3,4\}$ in row~2 and
$\{1,2\}$ in row~3. Both edges of $\rho_1$ join different
$G_1$-orbits. For $G_2$, only $\{5,6\}$ crosses.
The displayed supports give all claims about shared labels:
every label at least $3$ fixes the first five points.

Finally $G_3=\gen{A_5,B}$ is an actual product on disjoint
supports $D,E$. Its orbits are $D$ and all the $B$-orbits in
$E$, including fixed points. If $B$ is transitive, precisely
the transposition $(6\ b)$ in $\tau$ joins the two orbits;
all $\delta$-edges remain in $E$. Conversely a unique eligible
edge and transitivity of $G$ force exactly two omission orbits,
so $B$ is transitive. All prefix labels are confined to $D$
and all suffix labels to $E$, proving that this cut is unshared.
\end{proof}

\paragraph{Scope and source route.}
These are exact early-colour adapters for the exceptional-head
continuation of source Propositions~5.28--5.32, physical pages
36--45. The literal Table~2 data and the accepted critical-tail
and pure-overlap proofs remain paid. The new third-orbit argument
excludes a unique colour-$4$ crossing even when $B$ is intransitive.
It does not make $B$ natural alternating, eliminate a shared cut
of higher colour, or resolve the remaining full-source branches.
No new finite computation, formalization, original or question
credit is claimed.
\end{document}
